\documentclass[11pt]{amsart}

\usepackage[T1]{fontenc}
\usepackage{lmodern}
\usepackage{amsmath,amssymb,amsthm,mathtools}
\usepackage{enumitem}
\usepackage{booktabs}
\usepackage[margin=1.15in]{geometry}
\usepackage{xcolor}
\usepackage[colorlinks=true,linkcolor=blue!50!black,citecolor=blue!50!black,urlcolor=blue!50!black]{hyperref}

% ---------------------------------------------------------------------------
% Edit this macro once the repository is public.
\newcommand{\repourl}{\url{https://github.com/schlafly1/equivariant-embedding}}
% ---------------------------------------------------------------------------

\newtheorem{theorem}{Theorem}[section]
\newtheorem{proposition}[theorem]{Proposition}
\newtheorem{lemma}[theorem]{Lemma}
\newtheorem{corollary}[theorem]{Corollary}
\newtheorem*{maintheorem}{Main Theorem}
\theoremstyle{definition}
\newtheorem{definition}[theorem]{Definition}
\newtheorem{remark}[theorem]{Remark}

\newcommand{\R}{\mathbb{R}}
\newcommand{\E}{\mathbb{E}}
\newcommand{\Met}{\mathrm{Met}}
\newcommand{\Sym}{\operatorname{Sym}}
\newcommand{\Ric}{\operatorname{Ric}}
\newcommand{\tr}{\operatorname{tr}}
\newcommand{\range}{\operatorname{range}}
\newcommand{\inner}[2]{\langle #1, #2\rangle}
\newcommand{\cM}{\mathcal{M}}
\newcommand{\cB}{\mathcal{B}}
\newcommand{\cS}{\mathcal{S}}
\newcommand{\cT}{\mathcal{T}}
\newcommand{\cK}{\mathcal{K}}
\newcommand{\lean}[1]{\texttt{\detokenize{#1}}}

\title{Equivariant isometric embeddings via G\"unther's method}
\author{Roger Schlafly}
\date{\today}

\begin{document}

\begin{abstract}
Moore and Schlafly proved in 1980 that a compact Riemannian manifold on which a compact Lie group $G$
acts by isometries has an isometric embedding into a Euclidean space that is equivariant with respect
to an orthogonal representation of $G$. The original proof used the Nash--Moser implicit function
theorem with smoothing operators. We give a proof that replaces Nash--Moser by G\"unther's method,
which reduces the isometric embedding equation to a fixed point problem for a contraction on a
Banach space of H\"older maps. All of G\"unther's operators commute with the group action, so the
contraction preserves the closed subspace of equivariant maps, and the uniqueness of the small fixed
point makes the correction equivariant without any averaging of nonlinear solutions. We also
formalize the algebraic and functional-analytic core of the argument in Lean~4 with Mathlib:
the equivariant G\"unther fixed point theorem, the equivariance and coordinate independence of the
Nash--G\"unther pointwise right inverse, the freeness of the Veronese map, the Haar-averaged inner
product used in the approximation step, and the assembly of the main theorem from explicitly
stated analytic inputs. The formal proofs use only the standard axioms of Lean.
\end{abstract}

\subjclass[2020]{53C42 (primary); 57S15, 58J05, 35J60, 68V20 (secondary)}
\keywords{Isometric embedding, Nash embedding theorem, G\"unther's method, equivariant embedding,
compact Lie group, formalized mathematics, Lean}

\maketitle

\section{Introduction}

The Nash embedding theorem~\cite{Nash56} says that every Riemannian manifold has an isometric
embedding into a Euclidean space of sufficiently large dimension. In~\cite{MS80} J.~D.~Moore and
the author proved an equivariant version, extending the case of homogeneous spaces treated by Moore~\cite{Moore76}.

\begin{maintheorem}[Moore--Schlafly~\cite{MS80}]
Let $(M,g)$ be a compact Riemannian manifold and let $G$ be a compact Lie group acting smoothly on
$M$ by isometries. Then there are an orthogonal representation $\rho\colon G\to O(N)$ and a smooth
embedding $X\colon M\to\R^N$ such that
\[
X^*\inner{\cdot}{\cdot}_{\R^N} = g
\qquad\text{and}\qquad
X(\sigma p) = \rho(\sigma)\,X(p)\quad\text{for all }\sigma\in G,\ p\in M .
\]
\end{maintheorem}

The proof in~\cite{MS80} has two steps. Step~1 finds a ``perturbable'' (free) equivariant embedding
$X_0$ and shows, using the Nash--Moser implicit function theorem in the form of
Sergeraert~\cite{Sergeraert70}, that every invariant metric close to $F(X_0)$ is induced by an
equivariant embedding. Smoothing operators are made equivariant by composing with the Haar
projection. Step~2 shows that invariant metrics induced by equivariant maps are dense, using the
spectral theory of the Laplacian and Haar averaging.

In 1989 G\"unther~\cite{Gunther89,Gunther91} found a much simpler proof of the Nash embedding
theorem. He rewrote the linearized problem so that the nonlinear equation becomes a fixed point
problem
\begin{equation}\label{eq:fixed-point-intro}
v = \cM(h) + \cB(v,v)
\end{equation}
in which $\cB$ loses no derivatives. The ordinary contraction mapping principle on a H\"older space
then replaces the Nash--Moser iteration. Expositions are given by Yang~\cite{Yang98} and
Tao~\cite{Tao16}; both note that G\"unther's reduction to a torus is a convenience and that one can
work directly on $M$ with a background metric.

This paper does two things.

\begin{enumerate}[label=(\arabic*)]
\item \emph{A proof of the Main Theorem by G\"unther's method}
(Sections~\ref{sec:setup}--\ref{sec:proof}). We work intrinsically on $M$, using the Levi-Civita
connection of the $G$-invariant metric $g$ and the operator $(\Delta-1)^{-1}$. Every operator in
G\"unther's construction is then natural under isometries, so the map
$v\mapsto\cM(h)+\cB(v,v)$ commutes with the action of $G$ on maps $M\to\R^N$. Its unique small
fixed point is therefore equivariant (Theorem~\ref{thm:abstract-gunther}). No smoothing operators
are needed and no nonlinear solution is averaged. Step~2 is the same as in~\cite{MS80}.

\item \emph{A machine-checked formalization} of the algebraic and functional-analytic core of
this proof in Lean~4~\cite{Lean4} with Mathlib~\cite{Mathlib20} (Section~\ref{sec:lean}).
Mathlib does not yet contain H\"older spaces of sections over manifolds, the Laplace--Beltrami
operator, Schauder estimates or the Mostow--Palais theorem. We therefore prove the new
equivariant ingredients unconditionally, and we prove the Main Theorem conditionally: the
classical analytic inputs are explicit hypotheses, and none of them asserts the existence of an
equivariant solution of the nonlinear problem.
\end{enumerate}

Equivariant versions of G\"unther's construction for cocompact discrete group actions on
noncompact manifolds have been studied recently by Qiu~\cite{Qiu25}; see also
Burago--Qiu~\cite{BuragoQiu25}. As in~\cite{MS80}, no bound on $N$ depending only on $\dim M$ is
possible in the equivariant setting: \cite[\S6]{MS80} constructs left-invariant metrics on $S^3$
that need arbitrarily large $N$.

\subsection*{Conventions}
All manifolds are smooth, compact and without boundary (see Remark~\ref{rem:boundary} for
boundary), and $n=\dim M$. Repeated indices are summed and indices are raised with $g$. ``Smooth''
means $C^\infty$. We fix $\alpha\in(0,1)$.

\section{Setup}\label{sec:setup}

\subsection{Equivariant maps and induced metrics}
Let $(V,\rho)$ be a finite-dimensional orthogonal representation of $G$. A map $X\colon M\to V$
is \emph{equivariant} if $X(\sigma p)=\rho(\sigma)X(p)$. Its \emph{induced metric} is the symmetric
$2$-tensor
\[
F(X) = \inner{dX}{dX},\qquad F(X)_{ij} = \inner{\partial_i X}{\partial_j X}.
\]
If $X$ is equivariant then $F(X)$ is $G$-invariant, because $\rho(\sigma)$ is orthogonal. Two maps
$X_1\colon M\to V_1$ and $X_2\colon M\to V_2$ give $(X_1,X_2)\colon M\to V_1\oplus V_2$ with
\begin{equation}\label{eq:F-additive}
F(X_1,X_2) = F(X_1)+F(X_2),\qquad F(tX) = t^2F(X).
\end{equation}
If $X_1$ is an embedding, so is $(X_1,X_2)$. A $G$-invariant symmetric $2$-tensor is
\emph{realizable} if it equals $F(X)$ for some smooth equivariant map $X$ into some orthogonal
representation. By~\eqref{eq:F-additive} the realizable tensors form a convex cone.

\subsection{H\"older spaces and the group action}\label{subsec:holder}
Let $\nabla$ be the Levi-Civita connection of $g$; it acts on $V$-valued tensor fields
componentwise. For a $V$-valued tensor field $T$ we use the H\"older norm
\[
\|T\|_{k,\alpha} = \sum_{j=0}^{k}\sup_M|\nabla^jT| + [\nabla^kT]_\alpha,\qquad
[S]_\alpha = \sup_{0<d(x,y)<\iota}\frac{|S(x)-\Pi_{y\to x}S(y)|}{d(x,y)^\alpha},
\]
where $d$ is the Riemannian distance, $\iota$ is the injectivity radius and $\Pi_{y\to x}$ is
parallel transport along the unique minimizing geodesic. This norm is equivalent to the norm
defined with charts and a partition of unity. Let $C^{k,\alpha}(M;V)$ and
$C^{k,\alpha}(S^2T^*M)$ be the resulting Banach spaces.

The group $G$ acts on $V$-valued maps and on tensors by
\begin{equation}\label{eq:action}
(\sigma\cdot v)(x) = \rho(\sigma)\,v(\sigma^{-1}x),\qquad \sigma\cdot h = (\sigma^{-1})^*h .
\end{equation}
The fixed points of the first action are exactly the equivariant maps, and those of the second are
the invariant tensors.

\begin{lemma}\label{lem:isometric-action}
The actions~\eqref{eq:action} are linear and isometric for $\|\cdot\|_{k,\alpha}$.
\end{lemma}

\begin{proof}
An isometry $\sigma$ of $(M,g)$ preserves $\nabla$, the distance, minimizing geodesics, the
injectivity radius and parallel transport, so $\nabla^j(\sigma\cdot T)=\sigma\cdot\nabla^jT$ and
every term of the norm is preserved. The factor $\rho(\sigma)$ is orthogonal, so it preserves
pointwise norms.
\end{proof}

For an arbitrary equivalent norm, $\|v\|' = \sup_{\sigma\in G}\|\sigma\cdot v\|$ is an equivalent
norm for which $G$ acts isometrically. So the choice of norm is not essential.

\subsection{Elliptic facts}
Let $\Delta=\tr_g\nabla^2$ be the rough Laplacian on scalar or $V$-valued tensor fields. It is
nonpositive and self-adjoint, so $\Delta-1$ is invertible on each tensor bundle. We write $P_0$ and
$P_1$ for $(\Delta-1)^{-1}$ on symmetric $2$-tensors and on $1$-forms. We use these standard facts.

\begin{enumerate}[label=(E\arabic*)]
\item\label{E:schauder} (Schauder estimates) $P_0$ and $P_1$ are bounded from $C^{k,\alpha}$ to
$C^{k+2,\alpha}$ for every $k\ge0$; see e.g.~\cite[Ch.~6]{GT} for the local estimates and~\cite{Aubin} for compact manifolds.
\item\label{E:natural} (Naturality) $\Delta$, $P_0$, $P_1$ and $\Ric$ commute with isometries.
\item\label{E:products} (Products) Pointwise contractions define bounded bilinear maps
$C^{k,\alpha}\times C^{k,\alpha}\to C^{k,\alpha}$, and multiplication by a smooth tensor is bounded
on $C^{k,\alpha}$.
\item\label{E:bochner} (Commutators) For a $V$-valued function $f$,
$\Delta\nabla f = \nabla\Delta f + \Ric(\nabla f)$. For a $1$-form $w$, the commutator
$\cK(w) = \Sym([\Delta,\nabla]w)$ is a first-order operator whose coefficients are built from the
curvature tensor and its covariant derivative. In particular
$\cK\colon C^{k+1,\alpha}\to C^{k,\alpha}$ is bounded and commutes with isometries.
\end{enumerate}

\section{Free maps and the pointwise right inverse}\label{sec:right-inverse}

For a smooth map $u\colon M\to V$ and $x\in M$, define
\[
A_u(x)\colon T_xM\oplus S^2T_xM\to V,\qquad
A_u(x)(\xi,s) = du_x(\xi) + \nabla^2u_x(s),
\]
where the Hessian $\nabla^2u_x\in S^2T^*_xM\otimes V$ is contracted with $s$.

\begin{definition}
The map $u$ is \emph{free} if $A_u(x)$ is injective for every $x\in M$.
\end{definition}

In coordinates, $\partial_i\partial_ju = \nabla^2_{ij}u+\Gamma^k_{ij}\partial_ku$. So $u$ is free
if and only if the $n+\tfrac12n(n+1)$ vectors $\partial_iu$ and $\partial_i\partial_ju$
$(i\le j)$ are linearly independent. This is Nash's condition, called \emph{perturbable}
in~\cite{MS80}.

Give $T_xM\oplus S^2T_xM$ the inner product induced by $g$ and let $A_u(x)^\dagger$ be the
adjoint. After identifying vectors and covectors with $g$,
\begin{equation}\label{eq:adjoint}
A_u^\dagger y = \bigl(\inner{du}{y},\ \inner{\nabla^2u}{y}\bigr),\qquad
\text{i.e. } (A_u^\dagger y)_i = \inner{\partial_iu}{y},\quad (A_u^\dagger y)_{ij} = \inner{\nabla^2_{ij}u}{y}.
\end{equation}
For free $u$ define
\[
L_u(x) = A_u(x)\bigl(A_u(x)^\dagger A_u(x)\bigr)^{-1}\colon T^*_xM\oplus S^2T^*_xM\to V .
\]

\begin{proposition}[Lean: \lean{nash_gunther_right_inverse}]\label{prop:right-inverse}
Let $A\colon W\to V$ be an injective linear map of finite-dimensional real inner product spaces
and $b\in W$. Then $x=A(A^\dagger A)^{-1}b$ is the unique solution of $A^\dagger x=b$ in
$\range A$, and it has the least norm among all solutions of $A^\dagger x=b$.
\end{proposition}

\begin{proof}
$A^\dagger A$ is injective because $\inner{A^\dagger Aw}{w}=|Aw|^2$, hence invertible. If $x=Aw$
and $A^\dagger x=b$ then $w=(A^\dagger A)^{-1}b$. If $A^\dagger y=b$, then $y-x\in\ker
A^\dagger=(\range A)^\perp$, so $|y|^2=|x|^2+|y-x|^2$.
\end{proof}

So $L_u(x)$ is a right inverse of $A_u(x)^\dagger$, depends smoothly on $x$, and does not depend
on choices of coordinates. The next statement is the pointwise heart of equivariance in Step~1
of~\cite{MS80}.

\begin{proposition}[Lean: \lean{nash_gunther_right_inverse_equivariant}]\label{prop:right-inverse-eq}
Let $A,A'\colon W\to V$ with $A$ injective, let $R$ be a linear isometry of $V$ and $P$ a linear
automorphism of $W$, and suppose $A'=R\circ A\circ P$. If $x\in\range A$ solves $A^\dagger x=b$
and $x'\in\range A'$ solves $A'^\dagger x'=P^\dagger b$, then $x'=Rx$.
\end{proposition}

\begin{proof}
$Rx\in\range A'$ since $P$ is onto, and $A'^\dagger Rx = P^\dagger A^\dagger R^{-1}Rx=P^\dagger b$
since $R^\dagger=R^{-1}$. Apply the uniqueness in Proposition~\ref{prop:right-inverse}.
\end{proof}

\begin{corollary}\label{cor:L-equivariant}
Let $u$ be a free equivariant map and $b$ a $G$-invariant section of $T^*M\oplus S^2T^*M$ of class
$C^{k,\alpha}$. Then $x\mapsto L_u(x)b(x)$ is an equivariant map of class $C^{k,\alpha}$.
\end{corollary}

\begin{proof}
Differentiating $u\circ\sigma=\rho(\sigma)\circ u$ and using~\ref{E:natural} for $\nabla$ gives
$A_u(\sigma x)\circ Q = \rho(\sigma)\circ A_u(x)$, where $Q=d\sigma_x\oplus S^2d\sigma_x$ is an
isometry. Apply Proposition~\ref{prop:right-inverse-eq} with $A=A_u(x)$, $A'=A_u(\sigma x)$,
$R=\rho(\sigma)$ and $P=Q^{-1}$. Then $P^\dagger=Q$ and $Qb(x)=b(\sigma x)$ by invariance, so
$L_u(\sigma x)b(\sigma x)=\rho(\sigma)L_u(x)b(x)$. Regularity follows from~\ref{E:products}.
\end{proof}

\section{A free equivariant starting embedding}\label{sec:veronese}

By the theorem of Mostow and Palais~\cite{Mostow57,Palais57} there is an equivariant embedding
$e\colon M\to W$ into an orthogonal representation $W$. Give $S^2W\subset W\otimes W$ the induced
inner product, so that $S^2\rho$ is orthogonal, and put
\[
u_0 = (e,\ e\odot e)\colon M\to W\oplus S^2W,\qquad e\odot e = e\otimes e .
\]
This is the construction of~\cite[\S2]{MS80}; in an orthonormal basis it is
$(x^a,\ (x^a)^2,\ \sqrt2\,x^ax^b)_{a<b}$.

\begin{proposition}[Lean: \lean{veronese_map_free}]\label{prop:veronese}
$u_0$ is a free equivariant embedding.
\end{proposition}

\begin{proof}
Equivariance and injectivity are clear. Fix a point, write $x=e(p)$, $x_i=\partial_ie(p)$ (linearly
independent) and $x_{ij}=\partial_i\partial_je(p)$, and write $a\odot b=a\otimes b+b\otimes a$. By
the product rule,
\[
\partial_iu_0 = \bigl(x_i,\ x_i\odot x\bigr),\qquad
\partial_i\partial_ju_0 = \bigl(x_{ij},\ x_{ij}\odot x + x_i\odot x_j\bigr)
\]
(for $\partial_i\partial_j$ of $e\otimes e$ the last term is $x_i\otimes x_j+x_j\otimes x_i$).
Suppose $\sum c_i\partial_iu_0+\sum d_{ij}\partial_i\partial_ju_0=0$ with $d$ symmetric. The first
component gives $a:=\sum c_ix_i+\sum d_{ij}x_{ij}=0$. Evaluating the second component on
$\varphi\otimes\psi$ for linear functionals $\varphi,\psi$ gives
\[
\varphi(a)\psi(x)+\varphi(x)\psi(a)+2\sum_{i,j}d_{ij}\varphi(x_i)\psi(x_j)=0 .
\]
Choosing $\varphi,\psi$ dual to $x_k,x_l$ gives $d_{kl}=0$, and then $\sum c_ix_i=0$ gives
$c=0$.
\end{proof}

Since $tu_0$ is free for $t\ne0$, Proposition~\ref{prop:veronese} supplies free equivariant
embeddings with induced metric $t^2F(u_0)$.

\section{G\"unther's reformulation on a Riemannian manifold}\label{sec:gunther}

Let $u$ be a smooth free equivariant embedding into $V$. We look for $v\in C^{2,\alpha}(M;V)$ with
$F(u+v)=F(u)+h$, where $h$ is a given invariant symmetric $2$-tensor. Write
$(\Sym\nabla w)_{ij}=\tfrac12(\nabla_iw_j+\nabla_jw_i)$ for a $1$-form $w$.

\subsection{Two identities}
For $V$-valued $u$ and $v$,
\begin{equation}\label{eq:id1}
\inner{\nabla_iu}{\nabla_jv}+\inner{\nabla_ju}{\nabla_iv}
= 2\Sym\nabla\bigl(\inner{du}{v}\bigr)_{ij} - 2\inner{\nabla^2_{ij}u}{v},
\end{equation}
since $\nabla_i\inner{\nabla_ju}{v}=\inner{\nabla^2_{ij}u}{v}+\inner{\nabla_ju}{\nabla_iv}$ and the
Hessian is symmetric.

The quadratic term $a_{ij}=\inner{\nabla_iv}{\nabla_jv}$ is handled by G\"unther's trick. For
$V$-valued $v,w\in C^{2,\alpha}$ define the $1$-form and $2$-tensor
\begin{align*}
\sigma(v,w)_j &= \tfrac12\bigl(\inner{\nabla_jv}{\Delta w}+\inner{\nabla_jw}{\Delta v}\bigr),\\
\tilde\tau(v,w)_{ij} &= \Sym_{(ij),(vw)}\Bigl[-2\inner{\nabla^2_{ij}v}{\Delta w}
  + 2\inner{\nabla_i\nabla_kv}{\nabla_j\nabla^kw}\\
  &\qquad + 2\Ric_i^{\,k}\inner{\nabla_kv}{\nabla_jw} - \inner{\nabla_iv}{\nabla_jw}\Bigr],
\end{align*}
where $\Sym_{(ij),(vw)}$ symmetrizes in $i,j$ and in $v,w$. Both lie in $C^{0,\alpha}$
by~\ref{E:products}. Put
\begin{equation}\label{eq:S-T}
\cS(v,w) = P_1\,\sigma(v,w),\qquad
\cT(v,w) = P_0\bigl(\tilde\tau(v,w) - 2\,\cK(\cS(v,w))\bigr).
\end{equation}
By~\ref{E:schauder} and~\ref{E:bochner}, $\cS(v,w)\in C^{2,\alpha}(T^*M)$ and
$\cT(v,w)\in C^{2,\alpha}(S^2T^*M)$.

\begin{lemma}[G\"unther's identity]\label{lem:gunther-identity}
For $v\in C^{2,\alpha}(M;V)$,
\[
\inner{\nabla_iv}{\nabla_jv} = 2\Sym\nabla\bigl(\cS(v,v)\bigr)_{ij} + \cT(v,v)_{ij}.
\]
\end{lemma}

\begin{proof}
First let $v$ be smooth. By the Leibniz rule and~\ref{E:bochner}, applied to the $V$-valued
$1$-form $\nabla v$,
\[
\Delta a_{ij} = \inner{\nabla_i\Delta v}{\nabla_jv}+\inner{\nabla_iv}{\nabla_j\Delta v}
+\Ric_i^{\,k}\inner{\nabla_kv}{\nabla_jv}+\Ric_j^{\,k}\inner{\nabla_kv}{\nabla_iv}
+2\inner{\nabla_i\nabla_kv}{\nabla_j\nabla^kv}.
\]
Using $\inner{\nabla_i\Delta v}{\nabla_jv}=\nabla_i\inner{\nabla_jv}{\Delta v}
-\inner{\nabla^2_{ij}v}{\Delta v}$, this gives
$(\Delta-1)a = 2\Sym\nabla\sigma(v,v)+\tilde\tau(v,v)$. Let $S=\cS(v,v)=P_1\sigma(v,v)$. By the
definition of $\cK$, $(\Delta-1)\Sym\nabla S=\Sym\nabla\bigl((\Delta-1)S\bigr)+\cK(S)
=\Sym\nabla\sigma(v,v)+\cK(S)$. Hence
\[
(\Delta-1)\bigl(a-2\Sym\nabla S\bigr) = \tilde\tau(v,v)-2\cK(S),
\]
and applying $P_0$ gives the claim. For general $v\in C^{2,\alpha}$, choose $\beta\in(0,\alpha)$.
Smooth maps are dense in $C^{2,\alpha}$ for the $C^{2,\beta}$ norm, and both sides of the identity
are continuous from $C^{2,\beta}$ to $C^{0,\beta}$.
\end{proof}

\subsection{The fixed point equation}
Combining~\eqref{eq:id1} and Lemma~\ref{lem:gunther-identity},
\[
F(u+v)-F(u) = 2\Sym\nabla\bigl(\inner{du}{v}+\cS(v,v)\bigr)
-2\Bigl(\inner{\nabla^2u}{v}-\tfrac12\cT(v,v)\Bigr).
\]
So $F(u+v)=F(u)+h$ holds as soon as
\begin{equation}\label{eq:system}
\inner{du}{v} = -\cS(v,v),\qquad
\inner{\nabla^2u}{v} = \tfrac12\cT(v,v)-\tfrac12h .
\end{equation}
By~\eqref{eq:adjoint} this system is $A_u^\dagger v=\bigl(-\cS(v,v),\ \tfrac12\cT(v,v)-\tfrac12h\bigr)$.
Define
\[
\cM(h) = L_u\bigl(0,-\tfrac12h\bigr),\qquad
\cB(v,w) = L_u\bigl(-\cS(v,w),\ \tfrac12\cT(v,w)\bigr),
\]
where $L_u$ acts pointwise.

\begin{proposition}\label{prop:gunther-package}
Let $E=C^{2,\alpha}(M;V)$ and $\Met^{2,\alpha}=C^{2,\alpha}(S^2T^*M)$.
\begin{enumerate}[label=(\alph*)]
\item\label{pkg:bounded} $\cM\colon\Met^{2,\alpha}\to E$ is bounded linear and
$\cB\colon E\times E\to E$ is bounded bilinear.
\item\label{pkg:identity} If $v\in E$ satisfies $v=\cM(h)+\cB(v,v)$, then $F(u+v)=F(u)+h$.
\item\label{pkg:equivariant} If $h$ is $G$-invariant then $\cM(h)$ is $G$-fixed, and
$\cB(\sigma\cdot v,\sigma\cdot w)=\sigma\cdot\cB(v,w)$ for all $\sigma\in G$.
\end{enumerate}
\end{proposition}

\begin{proof}
\ref{pkg:bounded} follows from~\ref{E:schauder}--\ref{E:bochner}, because $L_u$ has smooth
coefficients. For~\ref{pkg:identity}, $A_u^\dagger L_u=\mathrm{id}$ gives
$A_u^\dagger v = (0,-\tfrac12h)+(-\cS(v,v),\tfrac12\cT(v,v))$, which is~\eqref{eq:system}. For
\ref{pkg:equivariant}, the operators $\nabla$, $\Delta$, $\Ric$, $P_0$, $P_1$ and $\cK$ and the
contractions with $g$ are natural under isometries~\ref{E:natural}, and the inner product of $V$
is $\rho$-invariant. So $\cS(\sigma\cdot v,\sigma\cdot w)=\sigma\cdot\cS(v,w)$, and likewise for
$\cT$. Corollary~\ref{cor:L-equivariant} (in its pointwise form) gives the equivariance of $L_u$.
\end{proof}

We also need two standard facts.

\begin{enumerate}[label=(F\arabic*)]
\item\label{F:regularity} (Regularity, G\"unther~\cite{Gunther89,Gunther91}) There is $r>0$ such
that if $h$ is smooth, $v\in E$, $\|v\|_{2,\alpha}<r$ and $v=\cM(h)+\cB(v,v)$, then $v$ is smooth.
The reason is that, after differentiating the equation $k$ times, the top-order derivatives of $v$
satisfy a \emph{linear} equation $w=f+\mathcal{L}_vw$, where $\|\mathcal{L}_v\|\le
C\|v\|_{2,\alpha}$ and $C$ does not depend on $k$.
\item\label{F:open} (Openness of embeddings) There is $\delta>0$ such that if $\|v\|_{2,\alpha}<\delta$
then $u+v$ is an embedding; see~\cite[Ch.~2]{Hirsch}.
\end{enumerate}

\section{The equivariant G\"unther perturbation theorem}\label{sec:perturbation}

\begin{theorem}[Lean: \lean{gunther_fixed_point_equivariant}]\label{thm:abstract-gunther}
Let $E$ be a real Banach space on which a monoid $G$ acts by linear isometries. Let
$\cB\colon E\times E\to E$ be bounded bilinear and $G$-equivariant, and let $m\in E$ be
$G$-fixed with $4\|\cB\|\,\|m\|<1$. Then $v=m+\cB(v,v)$ has a unique solution with
$\|v\|\le2\|m\|$, and this solution is $G$-fixed.
\end{theorem}

\begin{proof}
Let $T(v)=m+\cB(v,v)$ and $\bar B=\{\|v\|\le2\|m\|\}$. For $v,w\in\bar B$ we have
$\|T(v)\|\le\|m\|+4\|\cB\|\|m\|^2\le2\|m\|$, and
$\|T(v)-T(w)\| = \|\cB(v,v-w)+\cB(v-w,w)\|\le4\|\cB\|\|m\|\,\|v-w\|$. So $T$ is a contraction of
the complete set $\bar B$ and has a unique fixed point $v$ there. For $\sigma\in G$, the point
$\sigma\cdot v$ lies in $\bar B$ (the action is isometric) and
$T(\sigma\cdot v)=\sigma\cdot m+\sigma\cdot\cB(v,v)=\sigma\cdot T(v)=\sigma\cdot v$. By
uniqueness, $\sigma\cdot v=v$.
\end{proof}

Lean proves this for a general quadratic map $Q$ with $Q(0)=0$ and
$\|Q(v)-Q(w)\|\le C(\|v\|+\|w\|)\|v-w\|$, and deduces the bilinear case.

\begin{theorem}[Equivariant perturbation; Lean: \lean{equivariant_gunther_perturbation}]\label{thm:perturbation}
Let $u$ be a smooth free equivariant embedding. There is $\varepsilon>0$ such that for every
smooth $G$-invariant symmetric $2$-tensor $h$ with $\|h\|_{2,\alpha}<\varepsilon$ there is a smooth
equivariant embedding $X$ (for the same representation) with $F(X)=F(u)+h$.
\end{theorem}

\begin{proof}
Use the equivariant Banach space $E=C^{2,\alpha}(M;V)$ (Lemma~\ref{lem:isometric-action}) and the
operators of Proposition~\ref{prop:gunther-package}. Let $r,\delta$ be as in~\ref{F:regularity}
and~\ref{F:open}, let
\[
\eta=\min\Bigl(\tfrac12\min(\delta,r),\ \frac1{4\|\cB\|+1}\Bigr),\qquad \varepsilon=\frac{\eta}{\|\cM\|+1}.
\]
If $\|h\|_{2,\alpha}<\varepsilon$ then $m=\cM(h)$ is $G$-fixed with $\|m\|<\eta$, so
$4\|\cB\|\|m\|<1$. Theorem~\ref{thm:abstract-gunther} gives a $G$-fixed $v$ with
$v=m+\cB(v,v)$ and $\|v\|\le2\|m\|<\min(\delta,r)$. By~\ref{F:regularity}, $v$ is smooth; it is
equivariant since it is $G$-fixed; $X=u+v$ is an embedding by~\ref{F:open}; and $F(X)=F(u)+h$ by
Proposition~\ref{prop:gunther-package}\ref{pkg:identity}.
\end{proof}

\begin{remark}
Theorem~\ref{thm:perturbation} needs no averaging of solutions. Averaging is in fact useless
here: the average of two solutions of $F(X)=g$ is not a solution. Equivariance comes entirely from
the uniqueness of the small fixed point. Equivalently, the contraction preserves the closed
subspace $E^G$ of equivariant maps, and one may run the iteration inside $E^G$.
\end{remark}

\section{Approximation by realizable metrics}\label{sec:approx}

This step is that of~\cite[\S3]{MS80}; we include it for completeness and to point out the part
that is formalized.

\begin{proposition}\label{prop:approx}
Every smooth positive definite $G$-invariant metric $g'$ on $M$ is a limit, in every
$C^{k,\alpha}$ norm, of realizable metrics.
\end{proposition}

\begin{proof}
By the (non-equivariant) Nash embedding theorem, for instance in G\"unther's form, there are
smooth functions $y^1,\dots,y^q$ with $g'=\sum_i dy^i\otimes dy^i$. (The weaker statement that
$g'$ can be approximated by such sums, proved directly in~\cite[p.~58]{Nash56}, is enough.) Let
$\Delta_g$ be the Laplace--Beltrami operator of $g$. The partial sums of the eigenfunction
expansion of a smooth function converge in every Sobolev norm, hence in every $C^{k}$ norm. So
there are finite sums $f_i$ of eigenfunctions with
\[
\Bigl\|g'-\sum_i df_i\otimes df_i\Bigr\|_{k,\alpha}<\varepsilon .
\]
Each eigenspace of $\Delta_g$ is finite-dimensional and $G$-invariant ($G$ acts by isometries). So
the span $V_i$ of the translates $\sigma\cdot f_i$ is a finite-dimensional $G$-invariant space of
functions. Let $X_i\colon M\to V_i^*$ be the evaluation map $X_i(p)(f)=f(p)$; it is smooth and
equivariant for the dual representation. Let $\varphi_i\in V_i^{**}$ be evaluation at $f_i$, and
give $V_i^*$ the inner product
\begin{equation}\label{eq:averaged-ip}
\inner{\xi}{\eta}_i = \int_G\xi(\sigma\cdot f_i)\,\eta(\sigma\cdot f_i)\,d\sigma
= \int_G \varphi_i(\sigma^{-1}\xi)\,\varphi_i(\sigma^{-1}\eta)\,d\sigma,
\end{equation}
with $d\sigma$ normalized Haar measure. By Proposition~\ref{prop:avg-ip} below, this is a
$G$-invariant inner product, since the translates of $f_i$ span $V_i$. For tangent vectors $a,b$
at $p$,
\[
F(X_i)(a,b) = \int_G d(\sigma\cdot f_i)(a)\,d(\sigma\cdot f_i)(b)\,d\sigma
= \Pi\bigl(df_i\otimes df_i\bigr)(a,b),\qquad \Pi(\kappa)=\int_G\sigma^*\kappa\,d\sigma,
\]
where the Haar projection $\Pi$ is computed pointwise. Since $\Pi$ commutes with $\nabla$ and with
parallel transport, Lemma~\ref{lem:isometric-action} gives $\|\Pi\kappa\|_{k,\alpha}\le\|\kappa\|_{k,\alpha}$ (estimate each term of the norm inside the integral).
Also $\Pi g'=g'$. So $X=(X_1,\dots,X_q)$ is a smooth equivariant map with
\[
\|g'-F(X)\|_{k,\alpha}=\Bigl\|\Pi\Bigl(g'-\sum_idf_i\otimes df_i\Bigr)\Bigr\|_{k,\alpha}<\varepsilon .
\qedhere
\]
\end{proof}

\begin{proposition}[Lean: \lean{haar_averaged_inner_product}]\label{prop:avg-ip}
Let a compact group $G$ act linearly and continuously on a finite-dimensional real vector space
$W$. Let $\mu$ be a finite left-invariant measure on $G$ that is positive on nonempty open sets,
and let $\varphi\in W^*$ be such that the functionals $\varphi\circ\rho(\sigma)$ separate the
points of $W$. Then
$q(x,y)=\int_G\varphi(\rho(\sigma)^{-1}x)\,\varphi(\rho(\sigma)^{-1}y)\,d\mu(\sigma)$ is a
$G$-invariant inner product on $W$.
\end{proposition}

\begin{proof}
Bilinearity and symmetry are clear, and invariance follows from left-invariance of $\mu$. If
$x\ne0$, the continuous function $\sigma\mapsto\varphi(\rho(\sigma)^{-1}x)^2$ is nonnegative and
not identically zero, so its integral is positive.
\end{proof}

\begin{remark}
The construction can be made explicit. Decompose $L^2(M)$ (complex-valued) into irreducible
unitary representations of $G$ inside the eigenspaces, expand the approximating functions $z^i$ in
bases of these, and regroup the coefficients by isomorphism type. This gives an equivariant map $X$
with $\sum_{a,r}N_r^{1/2}X_{iaar}=z^i$ and $F(X)=\Pi\bigl(\sum_i dz^i\otimes dz^i\bigr)$;
see~\cite{Schlafly-notes}.
\end{remark}

\section{Proof of the Main Theorem}\label{sec:proof}

Let $g$ be the given invariant metric. Proposition~\ref{prop:veronese} gives a free equivariant
embedding $u_0$. Since $M$ is compact there is $c>0$ with $g-cF(u_0)$ positive definite. The map
$u_1=\sqrt c\,u_0$ is a free equivariant embedding with $F(u_1)=cF(u_0)$. Let $\varepsilon$ be
given by Theorem~\ref{thm:perturbation} for $u_1$. By Proposition~\ref{prop:approx} there is a
smooth equivariant map $X_2$ with
\[
\|h\|_{2,\alpha}<\varepsilon,\qquad h := g-cF(u_0)-F(X_2).
\]
The tensor $h$ is smooth and $G$-invariant. Theorem~\ref{thm:perturbation} gives a smooth
equivariant embedding $X_1$ with $F(X_1)=F(u_1)+h=g-F(X_2)$. Then $X=(X_1,X_2)$ is an equivariant
embedding into the orthogonal representation $V_1\oplus V_2$, and $F(X)=F(X_1)+F(X_2)=g$
by~\eqref{eq:F-additive}. \qed

This argument is formalized in Lean as \lean{equivariant_isometric_embedding_conditional}, with the
inputs listed in Table~\ref{tab:hypotheses} as hypotheses.

\begin{remark}\label{rem:boundary}
As in~\cite{MS80}, the case of manifolds with boundary reduces to the closed case: embed $M$
equivariantly in a sphere in an orthogonal representation and extend the invariant metric to an
invariant metric on the sphere by an invariant partition of unity. We do not treat the real-analytic
category, where~\cite[\S4]{MS80} uses the method of Greene and Jacobowitz.
\end{remark}

\section{The Lean formalization}\label{sec:lean}

The formalization is available at\\ \repourl.\\ It uses Lean~4 (toolchain \texttt{v4.35.0}) and
Mathlib (tag \texttt{v4.35.0}), and follows the layout of the Palomar registry: a short file
\texttt{Challenge.lean} with the advertised statements (proofs replaced by \texttt{sorry}), a file
\texttt{Solution.lean} with the same statements proved, and the proof development in the directory
\texttt{EquivEmbed/}. Every theorem in \texttt{Solution.lean} depends only on the axioms
\texttt{propext}, \texttt{Classical.choice} and \texttt{Quot.sound}; \texttt{Lean.ofReduceBool}
(\texttt{native\_decide}) is not used.

\subsection{Unconditional statements}

\begin{center}
\begin{tabular}{@{}lll@{}}
\toprule
Lean declaration & Mathematical content & Here\\
\midrule
\lean{gunther_fixed_point_equivariant} & equivariant G\"unther fixed point & Thm.~\ref{thm:abstract-gunther}\\
\lean{nash_gunther_right_inverse} & pointwise right inverse, least norm & Prop.~\ref{prop:right-inverse}\\
\lean{nash_gunther_right_inverse_equivariant} & equivariance, coordinates & Prop.~\ref{prop:right-inverse-eq}\\
\lean{veronese_map_free} & freeness of $(e,e\odot e)$ & Prop.~\ref{prop:veronese}\\
\lean{haar_averaged_inner_product} & averaged invariant inner product & Prop.~\ref{prop:avg-ip}\\
\bottomrule
\end{tabular}
\end{center}

These are stated with ordinary Mathlib notions: Banach spaces with an isometric action
(\texttt{IsIsometricSMul}), adjoints of linear maps between finite-dimensional inner product spaces,
linear functionals, Bochner integrals with respect to a left-invariant measure, and
\texttt{Representation}.

\subsection{The conditional main theorem}
Maps are elements of the structure \lean{EquivariantMap}: a function
$X\colon\texttt{Pt}\to\mathbb{R}^N$ on a type \texttt{Pt} with a $G$-action, bundled with a
homomorphism $\rho$ from $G$ to the linear isometries of $\mathbb{R}^N$ and a proof that
$X(\sigma p)=\rho(\sigma)X(p)$. Equivariance is thus enforced by the types, and the conclusions of
the theorems restate it. An abstract structure \lean{EmbeddingProblem} records the type
\texttt{Pt} (the points of $M$), a normed space \texttt{Met} (intended: invariant
$C^{2,\alpha}$ tensors), a subset \texttt{pos} (intended: positive definite tensors), the induced
metric $F$ from equivariant maps to \texttt{Met}, predicates for being an embedding (required to
imply injectivity) and being free, and the direct sum, with the identities~\eqref{eq:F-additive}
for the sum and for scaling. The theorem \lean{equivariant_isometric_embedding_conditional} states
that every element of \texttt{pos} equals $F(X)$ for an equivariant embedding $X$, under the four
hypotheses in Table~\ref{tab:hypotheses}.

The geometric meaning of \texttt{Met}, \texttt{pos}, $F$, freeness and being an embedding (beyond
injectivity) is not encoded in Lean. The theorem holds for every instantiation of
\lean{EmbeddingProblem}; its reading as a statement about Riemannian manifolds depends on an
instantiation for an actual compact Riemannian $G$-manifold, which is given by
Sections~\ref{sec:right-inverse}--\ref{sec:proof} but is not formalized.

\begin{table}[ht]
\caption{Hypotheses of the conditional main theorem.}\label{tab:hypotheses}
\begin{tabular}{@{}p{0.27\textwidth}p{0.42\textwidth}p{0.22\textwidth}@{}}
\toprule
Lean hypothesis & Content & Justification\\
\midrule
\lean{HasFreeEmbedding} & a free equivariant embedding exists & Mostow--Palais and Prop.~\ref{prop:veronese}\\
\lean{RealizableDense} & realizable metrics are dense in \texttt{pos} & Prop.~\ref{prop:approx}\\
\lean{ConeCondition} & $g-cF(X)\in\texttt{pos}$ for some $c>0$ (one $c$ suffices) & compactness of $M$\\
\lean{HasGuntherPackages} & at each free embedding: $E$, $\cM$, $\cB$ as in Prop.~\ref{prop:gunther-package}, \ref{F:regularity}, \ref{F:open} & Sections~\ref{sec:right-inverse}--\ref{sec:gunther}\\
\bottomrule
\end{tabular}
\end{table}

The structure \lean{GuntherPackage} asks for a Banach space $E$ with an isometric linear $G$-action,
a bounded linear $\cM$ with $G$-fixed values, a bounded bilinear equivariant $\cB$, a map sending
smooth $G$-fixed $v$ to a map ``$u+v$'', G\"unther's identity
(Proposition~\ref{prop:gunther-package}\ref{pkg:identity}), regularity of small solutions, and
openness of embeddings. These are non-equivariant analytic facts together with the naturality of
the operators; the existence of an equivariant solution is not assumed. It is produced in Lean by
Theorem~\ref{thm:abstract-gunther}, through \lean{equivariant_gunther_perturbation}.

\subsection{What is not formalized}
The objects of differential geometry (manifolds, H\"older spaces of sections, the Laplacian,
Schauder estimates, eigenfunction expansions, the Mostow--Palais theorem, the openness of
embeddings) are not formalized; they enter only through the hypotheses above. The link between the
abstract structure and an actual compact Riemannian $G$-manifold is therefore given by the paper,
not by Lean. Formalizing it would need substantial new Mathlib infrastructure, in particular
elliptic regularity on compact manifolds.

\subsection{Checking}
The project builds with \texttt{lake build}. The script \texttt{scripts/gen\_challenge.py}
generates \texttt{Challenge.lean} and \texttt{Solution.lean} from a single list of statements, and
copies the definitions of \texttt{EquivEmbed/Defs.lean} verbatim into \texttt{Challenge.lean}. We
checked that the fully elaborated types of all compared declarations agree in the two environments,
and that \texttt{\#print axioms} reports only the three standard axioms. The continuous integration
configuration runs Lean's \texttt{lake comparator} with the independent kernels shipped with the
toolchain.

\section*{Acknowledgements}
The 1980 theorem is joint work with John Douglas Moore. The Lean formalization and the drafting of
this note were done with the assistance of Claude (Anthropic), an AI system; the author takes
responsibility for the mathematical content and the statements of all results.

\begin{thebibliography}{99}

\bibitem{Aubin}
T.~Aubin, \emph{Some Nonlinear Problems in Riemannian Geometry}, Springer Monographs in
Mathematics, Springer, Berlin, 1998.

\bibitem{BuragoQiu25}
D.~Burago and H.~Qiu, \emph{On the geometry of Riemannian isometric embeddings}, preprint (2025),
arXiv:2507.23164.

\bibitem{GT}
D.~Gilbarg and N.~S.~Trudinger, \emph{Elliptic Partial Differential Equations of Second Order},
Classics in Mathematics, Springer, Berlin, 2001.

\bibitem{Gunther89}
M.~G\"unther, \emph{Zum Einbettungssatz von J.~Nash}, Math. Nachr. \textbf{144} (1989), 165--187.

\bibitem{Gunther91}
M.~G\"unther, \emph{Isometric embeddings of Riemannian manifolds}, Proceedings of the International
Congress of Mathematicians (Kyoto, 1990), Vol.~II, Math. Soc. Japan, Tokyo, 1991, 1137--1143.

\bibitem{Hirsch}
M.~W.~Hirsch, \emph{Differential Topology}, Graduate Texts in Mathematics \textbf{33}, Springer,
New York, 1976.

\bibitem{Lean4}
L.~de~Moura and S.~Ullrich, \emph{The Lean 4 theorem prover and programming language}, in:
Automated Deduction -- CADE 28, Lecture Notes in Computer Science \textbf{12699}, Springer, 2021,
625--635.

\bibitem{Mathlib20}
The mathlib Community, \emph{The Lean mathematical library}, in: Proceedings of the 9th ACM SIGPLAN
International Conference on Certified Programs and Proofs (CPP 2020), ACM, 2020, 367--381.

\bibitem{Moore76}
J.~D.~Moore, \emph{Equivariant embeddings of Riemannian homogeneous spaces}, Indiana Univ. Math. J.
\textbf{25} (1976), 271--279.

\bibitem{MS80}
J.~D.~Moore and R.~Schlafly, \emph{On equivariant isometric embeddings}, Math. Z. \textbf{173}
(1980), 119--133.

\bibitem{Mostow57}
G.~D.~Mostow, \emph{Equivariant embeddings in Euclidean space}, Ann. of Math. \textbf{65} (1957),
432--446.

\bibitem{Nash56}
J.~Nash, \emph{The imbedding problem for Riemannian manifolds}, Ann. of Math. \textbf{63} (1956),
20--63.

\bibitem{Palais57}
R.~S.~Palais, \emph{Imbedding of compact, differentiable transformation groups in orthogonal
representations}, J. Math. Mech. \textbf{6} (1957), 673--678.

\bibitem{Qiu25}
H.~Qiu, \emph{On equivariant isometric embeddings of Riemannian manifolds with symmetries},
preprint (2025), arXiv:2511.09931.

\bibitem{Schlafly-notes}
R.~Schlafly, \emph{Notes on constructing approximate embeddings} and \emph{Equivariant G\"unther
method}, notes accompanying~\cite{MS80}, \url{https://github.com/schlafly1/papers}.

\bibitem{Sergeraert70}
F.~Sergeraert, \emph{Une g\'en\'eralisation du th\'eor\`eme des fonctions implicites de Nash},
C. R. Acad. Sci. Paris S\'er. A (1970), 861--863.

\bibitem{Tao16}
T.~Tao, \emph{Notes on the Nash embedding theorem}, blog post, May 11, 2016,
\url{https://terrytao.wordpress.com/2016/05/11/notes-on-the-nash-embedding-theorem/}.

\bibitem{Yang98}
D.~Yang, \emph{Gunther's proof of Nash's isometric embedding theorem}, preprint (1998),
arXiv:math/9807169.

\end{thebibliography}

\end{document}
